% ===================================================================
% MT-04a: Minitest Schulfächer (Las asignaturas) - MUSTERLÖSUNG
% Spanisch Klasse 7
% ===================================================================

\documentclass[11pt, a4paper]{article}

\usepackage[utf8]{inputenc}
\usepackage[T1]{fontenc}
\usepackage[ngerman]{babel}
\usepackage{helvet}
\renewcommand{\familydefault}{\sfdefault}

\usepackage[margin=2cm]{geometry}
\usepackage{enumitem}
\usepackage{tabularx}
\usepackage{booktabs}

\usepackage{xcolor}
\usepackage{tcolorbox}
\tcbuselibrary{skins}

\usepackage{fancyhdr}
\usepackage{lastpage}

% FARBEN
\definecolor{spanisch}{HTML}{c41e3a}
\definecolor{grau}{HTML}{888888}
\definecolor{hellgrau}{HTML}{f5f5f5}
\definecolor{loesung}{HTML}{c62828}

\newcommand{\fachfarbe}{spanisch}
\newcommand{\thema}{Minitest: Las asignaturas -- MUSTERLÖSUNG}
\newcommand{\fach}{Spanisch}
\newcommand{\klasse}{7}
\newcommand{\maxpunkte}{26}
\newcommand{\zeit}{15 Minuten}
\newcommand{\datum}{\today}

\pagestyle{fancy}
\fancyhf{}
\renewcommand{\headrulewidth}{0.4pt}
\renewcommand{\footrulewidth}{0.4pt}

\lhead{\fach{} -- Klasse \klasse}
\chead{\textbf{MT-04a -- ML}}
\rhead{\datum}

\lfoot{\zeit}
\cfoot{}
\rfoot{Seite \thepage\ von \pageref{LastPage}}

\newcommand{\punktebox}[1]{\hfill\fbox{\small \_\_\_/#1}}
\newcommand{\luecke}[1]{\rule{#1}{0.4pt}}
\newcommand{\lsg}[1]{\textcolor{loesung}{\textbf{#1}}}

\newtcolorbox{infobox}{
    colback=hellgrau,
    colframe=\fachfarbe,
    boxrule=1pt,
    arc=0pt,
    left=8pt, right=8pt, top=8pt, bottom=8pt
}

\newenvironment{aufgabe}[1]{%
    \vspace{12pt}
    \noindent\textbf{\textcolor{\fachfarbe}{Aufgabe #1}}
    \vspace{6pt}
    \begin{list}{}{\leftmargin=0pt}
    \item[]
}{%
    \end{list}
}

\newcommand{\es}[1]{\textcolor{\fachfarbe}{\textbf{#1}}}

\begin{document}

% ===================================================================
% KOPFBEREICH
% ===================================================================

\begin{center}
    {\Large\bfseries\textcolor{loesung}{\thema}}
\end{center}

\vspace{5pt}

\noindent
\begin{tabularx}{\textwidth}{|X|X|X|}
\hline
\textbf{Name:} \lsg{LÖSUNG} &
\textbf{Klasse:} \klasse &
\textbf{Datum:} \luecke{2cm} \\
\hline
\end{tabularx}

\vspace{10pt}

\begin{infobox}
\textbf{Zeit:} \zeit{} | \textbf{Punkte:} \maxpunkte{} BE | \textbf{Hilfsmittel:} keine
\end{infobox}

\vspace{15pt}

% ===================================================================
% AUFGABE 1: Zuordnung Schulfächer
% ===================================================================

\begin{aufgabe}{1 -- Schulfächer zuordnen}
Ordne die spanischen Schulfächer der deutschen Bedeutung zu. \punktebox{6}
\end{aufgabe}

\begin{center}
\begin{tabular}{rcl}
\es{las matemáticas} & \lsg{B} & Englisch \\[0.8em]
\es{la historia} & \lsg{C} & Mathematik \\[0.8em]
\es{el inglés} & \lsg{A} & Geschichte \\[0.8em]
\es{la educación física} & \lsg{D} & Sport \\[0.8em]
\es{la música} & \lsg{E} & Musik \\[0.8em]
\es{el arte} & \lsg{F} & Kunst \\
\end{tabular}
\end{center}

\textcolor{loesung}{\textit{(6 BE: je 1 BE pro richtige Zuordnung)}}

\vspace{15pt}

% ===================================================================
% AUFGABE 2: Artikel einsetzen
% ===================================================================

\begin{aufgabe}{2 -- Artikel einsetzen}
Wähle den richtigen Artikel: \es{el}, \es{la} oder \es{las}. \punktebox{6}
\end{aufgabe}

\noindent
a) \lsg{el} español \hfill
b) \lsg{las} matemáticas \\[0.8em]
c) \lsg{la} biología \hfill
d) \lsg{el} alemán \\[0.8em]
e) \lsg{la} geografía \hfill
f) \lsg{la} informática \\

\textcolor{loesung}{\textit{(6 BE: je 1 BE pro richtigen Artikel)}}

\textcolor{loesung}{\textit{Hinweis: Sprachen (español, alemán, inglés) = el (maskulin). Fächer auf -ía/-ica = la (feminin). matemáticas = immer Plural!}}

\vspace{15pt}

% ===================================================================
% AUFGABE 3: Me gusta / Me gustan
% ===================================================================

\begin{aufgabe}{3 -- Me gusta oder me gustan?}
Wähle die richtige Form: \es{gusta} (Singular) oder \es{gustan} (Plural). \punktebox{8}
\end{aufgabe}

\noindent
a) Me \lsg{gusta} el español. \hfill
b) Me \lsg{gustan} las matemáticas. \\[0.8em]
c) No me \lsg{gusta} la historia. \hfill
d) Me \lsg{gusta} el arte. \\[0.8em]
e) No me \lsg{gusta} la biología. \hfill
f) Me \lsg{gusta} la música. \\[0.8em]
g) Me \lsg{gustan} el inglés y el español. \hfill
h) No me \lsg{gusta} la educación física. \\

\textcolor{loesung}{\textit{(8 BE: je 1 BE pro richtige Form)}}

\textcolor{loesung}{\textit{Hinweis: gustan bei las matemáticas (Plural) und bei zwei Fächern (el inglés y el español). Sonst gusta (Singular).}}

\vspace{15pt}

% ===================================================================
% AUFGABE 4: Übersetzung
% ===================================================================

\begin{aufgabe}{4 -- Übersetze ins Spanische}
Schreibe die spanische Übersetzung. \punktebox{6}
\end{aufgabe}

\noindent
a) Spanisch gefällt mir. \\[0.5em]
\lsg{Me gusta el español.} \\[1em]

\noindent
b) Mathe gefällt mir nicht. \\[0.5em]
\lsg{No me gustan las matemáticas.} \\[1em]

\noindent
c) Musik gefällt mir. \\[0.5em]
\lsg{Me gusta la música.} \\

\textcolor{loesung}{\textit{(6 BE: je 2 BE pro korrekten Satz -- 1 BE für Struktur, 1 BE für Vokabular)}}

% ===================================================================
% PUNKTEÜBERSICHT
% ===================================================================

\vspace{20pt}
\hrule
\vspace{10pt}

\noindent
\begin{tabularx}{\textwidth}{|X|c|c|c|c||c|}
\hline
\textbf{Aufgabe} & 1 & 2 & 3 & 4 & \textbf{Gesamt} \\
\hline
\textbf{max. Punkte} & 6 & 6 & 8 & 6 & \textbf{26} \\
\hline
\textbf{erreicht} & & & & & \\[0.8em]
\hline
\end{tabularx}

\vspace{15pt}

\noindent\textbf{Bewertungsschlüssel (14-NP):}

\begin{center}
\begin{tabular}{|c|c|c|c|c|c|c|c|c|c|c|c|c|c|c|}
\hline
\textbf{NP} & 14 & 13 & 12 & 11 & 10 & 9 & 8 & 7 & 6 & 5 & 4 & 3 & 2 & 1 \\
\hline
\textbf{\%} & 100 & 96 & 91 & 86 & 80 & 73 & 67 & 60 & 53 & 47 & 40 & 33 & 27 & 20 \\
\hline
\textbf{BE} & 26 & 25 & 24 & 22 & 21 & 19 & 17 & 16 & 14 & 12 & 10 & 9 & 7 & 5 \\
\hline
\end{tabular}
\end{center}

\end{document}
